\documentclass[12pt]{article}
\usepackage[utf8]{inputenc}
\usepackage[T1]{fontenc}
\usepackage{hyperref}
\title{Wauyai Language Resources for K-12 Curriculum Development}
\author{Compiled by Claude Sonnet 5.5 (claude-sonnet-5-5) for Andr\'as Kornai}
\date{}
\begin{document}
\maketitle

\section{Introduction}
Wauyai (ISO 639-3: wuy) is an Austronesian language spoken on Waigeo Island in Raja Ampat, Southwest Papua, Indonesia. Estimates are about 300 speakers (Global Recordings Network) to about 400 (Joshua Project). It is closely related to Kawe and Laganyan, and all three descend from a shared root called Ma'ya. The language is endangered and is not known to be taught in schools. Searches found very little material: no primers, textbooks, dictionaries or school programmes.

\section{Primary Educational Resources}
No primers, textbooks or school materials in or about Wauyai turned up in searches. Joshua Project reports no known written or published form of the language. A curriculum would therefore have to be built from scratch, working with the community.

\section{Linguistic and Reference Materials}
The University of Edinburgh holds morphosyntactic elicitation data on Wauyai (Ma'ya). It was recorded in February 2023 with two native speakers in Waimalata hamlet on Waigeo, as part of a British Academy fellowship on the Raja Ampat--South Halmahera branch of Austronesian. Topics include verb classes, serial verbs and numeral classifiers. See the record at \url{https://datashare.ed.ac.uk/handle/10283/8610?show=full}. This is research data, not teaching material.

\section{Religious and Community Resources}
Joshua Project has a people-group profile for the Wauyai: \url{https://joshuaproject.net/people_groups/19542}. It gives demographic and location information. It does not list any scripture or community organisation resources.

\section{Digital and Audio Resources}
The Global Recordings Network page for Wauyai (\url{https://globalrecordings.net/en/language/wuy}) reports no audio recordings available at the time of the search.

\section{Curriculum Development Context}
Very little documentation exists, and the speaker community is small. A native-language curriculum would first need community-led documentation: an agreed orthography, a word list, and recordings of stories and everyday speech. The Edinburgh data could be a starting point for collaboration with the researchers. Until then, Indonesian would be the realistic language of instruction, with Wauyai introduced as a heritage subject.

\end{document}